\chapter{Boolean Analysis — Span}
%%% generated by the blueprint pipeline from TCSlib.BooleanAnalysis.GaussianSpace.Hermite.Span
\section{Overview}

Declarations of the Lean module \texttt{TCSlib.BooleanAnalysis.GaussianSpace.Hermite.Span}.

\section{Declarations}

\begin{definition}[Normalized Hermite polynomial sequence]
\lean{BooleanAnalysis.GaussianSpace.normalizedHermiteSequence}
\statementsource{boolean-ch11-gaussian-invariance}{pdf-8b2c50cf9621-p012-b001}
\label{BooleanAnalysis.GaussianSpace.normalizedHermiteSequence}
\leanok
\uses{BooleanAnalysis.GaussianSpace.normalizedHermitePolynomial, BooleanAnalysis.GaussianSpace.normalizedHermitePolynomial_degree}
The normalized Hermite polynomials, indexed by the natural numbers and with the degree
of the $k$th polynomial equal to $k$, define a sequence of polynomials over $\mathbb R$.
\end{definition}

\begin{definition}[Normalized Hermite polynomial basis]
\lean{BooleanAnalysis.GaussianSpace.normalizedHermiteBasis}
\statementsource{boolean-ch11-gaussian-invariance}{pdf-8b2c50cf9621-p012-b001}
\label{BooleanAnalysis.GaussianSpace.normalizedHermiteBasis}
\leanok
\uses{BooleanAnalysis.GaussianSpace.normalizedHermitePolynomial_ne_zero, BooleanAnalysis.GaussianSpace.normalizedHermiteSequence}
The normalized Hermite sequence is a basis of the real vector space of real polynomials,
indexed by the natural numbers.
\end{definition}

\begin{theorem}[Action of the normalized Hermite basis]
\lean{BooleanAnalysis.GaussianSpace.normalizedHermiteBasis_apply}
\label{BooleanAnalysis.GaussianSpace.normalizedHermiteBasis_apply}
\leanok
\uses{BooleanAnalysis.GaussianSpace.normalizedHermiteBasis, BooleanAnalysis.GaussianSpace.normalizedHermitePolynomial, BooleanAnalysis.GaussianSpace.normalizedHermiteSequence}
\difficulty{1}
For each nonnegative integer $k$, let $H_k$ be the polynomial obtained by multiplying
the $k$th Hermite polynomial, with its integer coefficients viewed as real coefficients,
by $(k!)^{-1/2}$. The $k$th element of the basis of $\mathbb R[X]$ formed from these
normalized Hermite polynomials is exactly $H_k$.
\end{theorem}

\begin{lemma}[Coordinate polynomial products span the Hermite family]
\lean{BooleanAnalysis.GaussianSpace.coordinatePolynomialProduct_mem_span_hermite}
\label{BooleanAnalysis.GaussianSpace.coordinatePolynomialProduct_mem_span_hermite}
\leanok
\uses{BooleanAnalysis.GaussianSpace.multivariateHermitePolynomial, BooleanAnalysis.GaussianSpace.normalizedHermiteBasis, BooleanAnalysis.GaussianSpace.normalizedHermiteBasis_apply}
\difficulty{6}
Let $n$ be a nonnegative integer and, for each coordinate $i\in\{0,\ldots,n-1\}$, let
$p_i$ be a real polynomial. The product over coordinates of the polynomial $p_i$ in
coordinate $i$ belongs to the real linear span of all multivariate Hermite polynomials
on $\mathbb R^n$, formed as products of normalized one-variable Hermite polynomials.
\end{lemma}

\begin{theorem}[Spanning by multivariate Hermite polynomials]
\lean{BooleanAnalysis.GaussianSpace.span_multivariateHermitePolynomial_eq_top}
\label{BooleanAnalysis.GaussianSpace.span_multivariateHermitePolynomial_eq_top}
\leanok
\uses{BooleanAnalysis.GaussianSpace.coordinatePolynomialProduct_mem_span_hermite, BooleanAnalysis.GaussianSpace.multivariateHermitePolynomial}
\difficulty{4}
For every nonnegative integer $n$, the normalized multivariate Hermite polynomials
indexed by multi-indices $\alpha\in\mathbb N^n$ span the real vector space of
polynomials in $n$ variables.
\end{theorem}
